% 최종 축약본. chapter4.md와 같은 본문에서 생성.
\subsection{자동 답변 품질 및 자원 사용 특성}
\label{sec:ch4_answer_resources}

Jev는 답변의 자동 평가에서도 평균 점수가 높았다(표 2). LF-Eval은 2.58점, Token F1은 1.12점 증가하였다. Token F1은 정답과 생성 답변의 공백 단위 표현 중복을 측정하고, LF-Eval은 정답 답변과 정답 근거를 기준으로 LLM이 답변을 평가한다. LF-Eval은 답변이 정답의 법률적 결론이나 근거를 빠뜨리거나 모순되게 서술했는지도 평가한다. 따라서 두 점수의 증가는 각각 표현의 일치와 정답 기준 평가의 개선을 뜻한다. 두 평가에는 공식 코드와 같은 앞 800자 평가 규칙을 적용하였다.

다만 Token F1의 별도 양측 검정 p값은 약 0.0525였다. 표의 신뢰구간과 이 검정은 산출 방식이 달라 결과가 일치하지 않을 수 있으므로, 모든 자동 지표가 통계적으로 유의하게 개선되었다고 해석하지 않았다. 또한 LF-Eval은 실제 선택 근거가 답변을 지지하는지를 직접 평가하지 않는다. 사람 평가는 수행되지 않아 전문가가 판단한 정확성과 근거 지지 여부는 아직 확인되지 않았다.

답변 생성에 제공한 근거의 평균 길이는 CE 203.67토큰, Jev 215.60토큰으로 늘었다. 질문·지시문·예시를 포함한 전체 입력도 평균 2302.67토큰에서 2314.58토큰으로 늘었고, 출력은 72.16토큰과 73.98토큰이었다. 사용량 토큰은 앞서 Token F1에서 사용한 공백 단위와 구별된다. 근거 길이의 중앙값은 각각 164개와 170.5개였고, 질문별 길이 차이의 중앙값은 0개였다. 평균 길이의 증가가 모든 질문에서 발생한 것은 아니다. 선택 목록의 길이는 문항별로 같았으므로, 이는 근거 수를 줄인 결과가 아니라 선택 내용이 달라진 결과이다.

따라서 이번 비교에서 입력량 절감은 관찰되지 않았다. 저장된 생성 응답을 재사용한 질문은 CE 72개, Jev 122개였다. 이 경우에도 동일 입력에 필요한 토큰량을 집계했으므로 신규 청구량을 뜻하지 않는다. 선택·점수화·평가·재시도까지 포함한 전체 서비스 비용은 산출하지 않았다. 실행 시간에는 저장 응답 재사용과 요청 제한 대기가 포함되어 있어 공정한 응답 지연 비교에도 사용할 수 없다. 또한 근거 길이의 증가가 품질 개선의 원인인지는 이 결과로 판단할 수 없다.
